
\begin{table}[ht]
	\centering
	\caption{مقایسه پژوهش حاضر با شکاف‌های باز پژوهش‌های پایه}
	\label{tab:open-issues-comparison}
	\resizebox{\textwidth}{!}{%
		\begin{tabular}{p{4.1cm}p{4.5cm}p{5.2cm}}
			\toprule
			\textbf{پژوهش پایه} &
			\textbf{مسئله یا مسیر آینده مطرح‌شده} &
			\textbf{پاسخ پژوهش حاضر} \\
			\midrule
			
			Coulibaly \& Hamlich (2026)
			\cite{coulibaly2026madtdt}
			&
			هماهنگ‌سازی سلسله‌مراتبی دوقلوهای دیجیتال، برآورد
			\lr{RUL}
			و استفاده از
			\lr{Explainable AI}
			در سامانه‌های چندعاملی
			&
			پیاده‌سازی معماری چندعاملی شامل دوقلوی دیجیتال، پیش‌بینی
			\lr{RUL}
			و عامل کاربردی؛ ارائه خروجی پیش‌بینی و اطلاعات مؤثر بر آن به کاربر و ایجاد بستر تعامل و تصمیم‌گیری. \\
			
			مطالعات دوقلوهای دیجیتال چندعاملی
			&
			افزایش استقلال، توزیع وظایف و قابلیت توسعه عامل‌ها
			&
			تفکیک خدمات در قالب عامل‌های مستقل و اتصال آن‌ها از طریق معماری لایه‌ای؛ امکان جایگزینی یا توسعه هر عامل بدون بازطراحی کل سامانه. \\
			
			مطالعات
			\lr{LLM--DT}
			&
			ایجاد واسط طبیعی و افزایش قابلیت تعامل انسان با دوقلوی دیجیتال
			&
			اتصال عامل مدل زبانی به وضعیت و خدمات دوقلوی دیجیتال؛ امکان پرس‌وجو، دریافت توضیح و ارسال فرمان با زبان طبیعی. \\
			
			مطالعات پیش‌بینی
			\lr{CMAPSS}
			&
			بهبود دقت پیش‌بینی
			\lr{RUL}
			به‌عنوان مسئله اصلی
			&
			استفاده از
			\lr{LSTM}
			به‌عنوان مؤلفه پیش‌بینی یک سامانه جامع؛ توجه به هزینه نامتقارن خطا و کاهش هدفمند کم‌برآورد/برآورد ناایمن. \\
			
			\textbf{پژوهش حاضر}
			&
			اتصال این مسیرهای پژوهشی در یک نمونه عملیاتی
			&
			\textbf{دوقلوی دیجیتال چندعاملی + LSTM + پایش شبه‌بی‌درنگ + تعامل زبانی + توضیح خروجی + معماری ماژولار و قابل تعمیم} \\
			
			\bottomrule
	\end{tabular}}
\end{table}
